\section{Arm estimates and multiscale address entropy}
\label{arm:section}

We now bound the number of contour neighborhoods needed to describe
adaptively chosen finite chord diagrams.  The fresh-block estimates of
Section~\ref{arm:inputs} give the probability cost of a prescribed
diagram.  A nested interval code pays for the choices of its locations.
For the fixed parameter,
\[
 \beta:=\frac12+\frac2{\gamma^2}>1.
\]
One unit of logarithmic address resolution costs one unit of entropy,
whereas one unit of monochrome rank costs at least $\beta$ units of
probability.  This strict surplus will bound total address occupancy and
force an acyclic component at a suitable scale.

We first replace arbitrary scalar contact graphs by chronological
comparisons with a total oscillation bound.  We then construct and count
the address code, apply the fresh-slot estimates, and allocate the
resulting entropy budget to a uniformly sampled scale.  All of these
steps concern the raw inventory word; they use no endpoint metric bound.

\subsection{Sorting scalar contacts with a total oscillation bound}

We first give a deterministic statement.  In particular, its constant does
not depend on the degrees of the contact graph.  A \emph{minimum chord} of
a continuous function $f$ is a pair $(u,v)$ such that
\[
 f(u)=f(v)=\min_{[u\wedge v,u\vee v]}f.
\]
All graphs in the next Lemma are finite.  Parallel edges and repeated chord
endpoints are allowed.

\begin{lemma}[Sorted scalar contacts]\label{arm:sorted-scalar-contacts}
Let $f$ be continuous on a compact interval containing pairwise disjoint
nonempty compact intervals $I_1,\ldots,I_m$, ordered from left to right.
Choose $t_v\in I_v$ and write
\[
 h_v=f(t_v),\qquad \omega_v=\max_{I_v}f-\min_{I_v}f.
\]
Let $G$ be a graph on these intervals, each of whose edges is witnessed by
a minimum chord with one endpoint in each incident interval.  Let $C$ be
a connected component of $G$, and list its vertices in increasing order
as $v_1,\ldots,v_k$.  For $1\le i<k$ set
\begin{equation}\label{arm:sorted-contact-errors}
 \mu_i=\min_{[t_{v_i},t_{v_{i+1}}]} f,
 \qquad E_i=h_{v_i}+h_{v_{i+1}}-2\mu_i.
\end{equation}
Then $E_i\ge0$,
\begin{equation}\label{arm:sorted-contact-walls}
 |h_{v_{i+1}}-h_{v_i}|\le E_i,
 \qquad
 \inf_{[t_{v_i},t_{v_{i+1}}]}f
 \ge \max\{h_{v_i},h_{v_{i+1}}\}-E_i,
\end{equation}
and
\begin{equation}\label{arm:sorted-contact-total}
 \sum_{i=1}^{k-1}E_i\le 6\sum_{v\in C}\omega_v.
\end{equation}
Thus the wall bound also holds on the gap between the two intervals, or
on any part of that gap obtained by deleting other intervals.  The same
bound can be summed over all components of $G$.
\end{lemma}

\begin{proof}
Discard any loops of $G$, which do not change its connected components.
We construct a finite metric tree directly from the minima of $f$.  This
also explains the contour order used in the proof.

Let $[A,B]$ contain all the intervals, and let
$m_*=\min_{[A,B]}f$.  For each $I_v=[a_v,b_v]$, choose a point $r_v$ at
which $f$ attains its minimum on $I_v$.  Let $\mathcal T$ be the finite set
consisting of $A,B$, all $a_v,b_v,r_v,t_v$, all chord endpoints, and a
point attaining $m_*$.  Coincident marked times occur only once in this
set.  Write
\[
 M(s,t)=\min_{[s\wedge t,s\vee t]}f\qquad(s,t\in\mathcal T).
\]
For each $t\in\mathcal T$ take a vertical segment
$\{t\}\times[m_*,f(t)]$.  Identify $(s,z)$ with $(t,z)$ exactly when
$z\le M(s,t)$.  This is an equivalence relation: the only nontrivial
check follows from
\[
 M(s,u)\ge\min\{M(s,t),M(t,u)\},
\]
because the interval between $s$ and $u$ is contained in the union of the
intervals between $s,t$ and between $t,u$.

The quotient, with distance measured along the vertical segments, is a
finite rooted metric tree $T$.  To verify this explicitly, adjoin the
marked times in increasing order.  The segment for a new time $t_j$
shares with the previously constructed tree precisely the initial
segment up to height $M(t_{j-1},t_j)$: indeed
\[
 \max_{i<j}M(t_i,t_j)=M(t_{j-1},t_j).
\]
Its remaining segment is therefore attached at a single point.  This
induction constructs a tree, allowing a new segment to have length zero.
Declare all marked tips and all attachment points to be vertices of $T$.
All its segments start at the common root of height $m_*$.  Denote the
top of the segment for $t$ by $\pi(t)$.  The common part of the two root
paths for $\pi(s),\pi(t)$ ends at height $M(s,t)$, so
\begin{equation}\label{arm:finite-contour-tree-distance}
 d_T(\pi(s),\pi(t))=f(s)+f(t)-2M(s,t).
\end{equation}
In particular, the endpoints of every minimum chord have the same image
in $T$.

For $I_v=[a_v,b_v]$, put $m_v=f(r_v)$ and define the subtree
\[
 L_v=[\pi(a_v),\pi(b_v)]_T
       \mathbin{\cup}[\pi(r_v),\pi(t_v)]_T.
\]
Here and below $[x,y]_T$ denotes the unique tree arc between $x$ and $y$.
The point $\pi(r_v)$ lies on the first arc: it is the common ancestor of
$\pi(a_v)$ and $\pi(b_v)$ at height $m_v$.  It is also an ancestor of
$\pi(t_v)$.  Thus $L_v$ is connected.  If $\ell$ denotes the length of a
finite metric tree, then
\begin{align}\label{arm:box-skeleton-length}
 \ell(L_v)
 &\le f(a_v)+f(b_v)-2m_v+f(t_v)-m_v\\
 &\le 3\omega_v.\notag
\end{align}
Overlapping arcs only reduce the first quantity.

Every chord contact in $I_v$ with a time outside $I_v$ belongs to
$[\pi(a_v),\pi(b_v)]_T$.  To see this, let $u\in I_v$ be the contact
and let $w<a_v$ be the other endpoint.  The chord condition implies
\[
 M(a_v,u)=f(u),\qquad m_v\le f(u)\le f(a_v).
\]
Hence $\pi(u)$ is the point of height $f(u)$ on the root path to
$\pi(a_v)$, between $\pi(r_v)$ and $\pi(a_v)$.  If instead $w>b_v$,
the same argument uses $b_v$.  This proves the claim, including contacts
at interval endpoints.

For each edge of $G$ joining $v$ to $w$, its common chord image therefore
belongs to $L_v\cap L_w$.  Graph connectivity shows that
\[
 U_C=\bigcup_{v\in C}L_v
\]
is connected.  A connected subset of a tree contains the unique arc
between any two of its points: removing an omitted point of that arc
would separate the two points.  Consequently $U_C$ is a subtree and
contains the finite hull $H$ of
$\pi(t_{v_1}),\ldots,\pi(t_{v_k})$.  Moreover,
\begin{equation}\label{arm:connected-box-skeleton-length}
 \ell(H)\le\ell(U_C)\le3\sum_{v\in C}\omega_v.
\end{equation}

We next justify the contour tour bound for the chronologically ordered
representatives.  Subdivide $H$ at every vertex of $T$ lying in $H$,
including all representative images; there are finitely many resulting
edges.  For a point $x$ in the interior of such an edge, let $\lambda$ be
its absolute height.  Among the components of $T\setminus\{x\}$, call the one
not containing the root the distal component.  The representative
indices whose images belong to this component form a consecutive block.
Indeed, if the images at times $s<u$ belong to it, then their common
ancestor has absolute height strictly greater than $\lambda$.
Thus $M(s,u)>\lambda$.  For any marked time $t\in[s,u]$,
\[
 M(s,t)\ge M(s,u)>\lambda,
\]
which puts $\pi(t)$ in the same distal component.  The strict inequality
is valid because $x$ is in an edge interior and all marked images are
vertices of the finite tree after subdivision.

The sequence of consecutive representative arcs can therefore cross any
edge of $H$ at most twice: it can enter and leave the consecutive block
of indices on its distal side only once each.  Summing edge lengths
gives
\begin{equation}\label{arm:chronological-tree-tour}
 \sum_{i=1}^{k-1}
 d_T\bigl(\pi(t_{v_i}),\pi(t_{v_{i+1}})\bigr)
 \le2\ell(H).
\end{equation}
This is the outer-face tour inequality, with its ordering verified
directly from interval minima.  It also constructs the relevant planar
order: root $H$ at its point nearest the root of $T$, order each child
subtree by its consecutive block of marked indices, and put repeated
marks of an internal vertex in the intervening corners.  The boundary
tour of this ordered tree encounters the marks chronologically and uses
each positive-length edge twice.  The preceding edge-count argument
does not require choosing any of those corners.

Equations~\eqref{arm:finite-contour-tree-distance},
\eqref{arm:connected-box-skeleton-length}, and
\eqref{arm:chronological-tree-tour} give
Equation~\eqref{arm:sorted-contact-total}.  Finally,
\[
 E_i=(h_{v_i}-\mu_i)+(h_{v_{i+1}}-\mu_i)
\]
is a sum of two nonnegative numbers.  It dominates their difference in
absolute value and each summand separately, proving
Equation~\eqref{arm:sorted-contact-walls}.  A singleton component has an
empty sum and needs no separate argument.
\end{proof}

\begin{corollary}[Discarding large comparison errors]
\label{arm:sorted-contact-discard}
Under the hypotheses of Lemma~\ref{arm:sorted-scalar-contacts}, for every
$\eta>0$ the number of consecutive comparisons with $E_i>\eta$, summed
over all components, is at most
\[
 \frac{6}{\eta}\sum_{v=1}^m\omega_v.
\]
All remaining comparisons have both representative height error and
wall error at most $\eta$.
\end{corollary}

\begin{proof}
Sum $\eta\mathbf 1_{\{E_i>\eta\}}\le E_i$ and apply
Equation~\eqref{arm:sorted-contact-total}.
\end{proof}

\begin{remark}\label{arm:sorted-contact-adaptedness}
The errors $E_i$ may depend on the whole contour.  The Lemma is a
deterministic implication, not a claim that these errors are measurable
with respect to the past of a sampling slot.  In a probabilistic use one
first keeps only comparisons whose error is at most a prescribed
deterministic tolerance.  The consequent wall may then be stated using
the earlier representative height and that tolerance.  Any conditioning
or enumeration needed to select the retained comparisons must still be
justified in the probability argument.
\end{remark}

\subsection{A noncrossing forest witnessing the contacts}

The preceding proof works for the original graph components, without
assuming that their witnesses form a noncrossing family.  For counting
purposes, the following separate construction gives a simultaneous
noncrossing witness, with possible mergers of components.

\begin{lemma}[Uncrossing and contraction]
\label{arm:noncrossing-contact-forest}
In the setting of Lemma~\ref{arm:sorted-scalar-contacts}, one may add
minimum chords whose endpoints are among the original chord endpoints
and then retain an outerplanar forest on the interval vertices with the
following properties.
\begin{enumerate}[label=\textup{(\roman*)}]
 \item The forest has one tree in each component of the enlarged chord
 graph.  Every original graph component is contained in one of these
 components.
 \item Every forest edge is witnessed by a genuine minimum chord.
 \item The vertices occur on the outer boundary in their time order.
 There are at most $32^m$ possible simple forests of this form on $m$
 prescribed ordered vertices.
\end{enumerate}
If the construction is applied to the endpoints belonging to a single
original connected component, it produces a noncrossing spanning tree
on exactly that component's vertices.  Forests obtained separately in
this way need not be noncrossing with one another.
\end{lemma}

\begin{proof}
On the finite set of original chord endpoints, put $s\sim t$ when
$f(s)=f(t)=M(s,t)$.  This is an equivalence relation, either directly
from the interval-minimum rule or from
Equation~\eqref{arm:finite-contour-tree-distance}.  For each class, add
the minimum chords between its consecutive times.  They connect all
times of that class, so they connect the endpoints of every original
chord.

These new chords can be drawn without crossings in the upper half-plane.
In fact, suppose two minimum chords have strictly interleaving endpoints
$a<b<c<d$, with pairs $(a,c)$ and $(b,d)$ at heights $h$ and $h'$.  Since
$b\in[a,c]$ and $c\in[b,d]$, the minimum conditions give respectively
$h'\ge h$ and $h\ge h'$.  Both heights agree, and all four times belong
to the same equivalence class.  Consecutive pairs of one class do not
interleave.  Chords from different classes therefore do not interleave
either.  Open semicircular arcs with nested or disjoint time intervals
as diameters do not intersect.  Shared
endpoints belong to the same class and cause no crossing.

Contract each interval $I_v$ in the boundary line to a single point.
This preserves an embedding with all vertices on the outer boundary.
One explicit way to see this is to choose a continuous proper
nondecreasing function $\varphi$ on the line that is constant precisely
on the nondegenerate $I_v$ and strictly increasing elsewhere.  For
$y>0$, the map
\[
 (x,y)\longmapsto
 \left(\frac{\varphi(x)+yx}{1+y},y\right)
\]
is a homeomorphism of the open upper half-plane, extends continuously
to the boundary, and there makes exactly the desired contractions.
Thus distinct chord interiors remain disjoint.  Discard loops and keep
one edge from each collection of parallel edges.  In every resulting
component choose a spanning tree, retaining its plane embedding.  The
earlier connectivity observation proves the first two assertions.
Using only endpoints from a single original connected component gives
the final assertion as well.

For completeness, the number of simple noncrossing forests on a fixed
ordered vertex set is exponential in its size.  Place the vertices on a
convex $m$-gon in that cyclic order.  Any noncrossing forest extends to a
triangulation after missing polygon sides are added and the remaining
faces are triangulated.  For $m\ge3$ there are fewer than $4^m$
triangulations, by the Catalan recursion, and each has fewer than $3m$
edges including the polygon boundary.  Its edge subsets number at most
$2^{3m}$.  Their product is at most $32^m$.  The cases $m=1,2$ satisfy
the same bound directly.
\end{proof}

\begin{remark}\label{arm:contact-mergers}
A simultaneous noncrossing forest can require merging original
components.  For example, two specified chords with interleaving
endpoints have a common contour height and allow additional true
chords between their endpoints, even if the specified graph initially
placed the two pairs in different components.  Such mergers are
explicit in Lemma~\ref{arm:noncrossing-contact-forest}; they must not be
silently treated as preservation of the original partition.  Adding
these true constraints cannot decrease the rank of the monochrome
relation.  Lemma~\ref{arm:sorted-scalar-contacts} itself needs no mergers.
\end{remark}


\subsection{Sparse projected addresses and disjoint sampling slots}

We now give a finite formulation of the address count.  Throughout this
subsection, $\mathbb P$ denotes the raw bilateral inventory-word law.
Additional randomness used to select finite witness lists may be included,
provided that it does not change the word's marginal law.  A subsequent
selection law $\mathbb Q_N$ will be required to satisfy
\begin{equation}\label{arm:address-tilt}
 \frac{\mathrm d\mathbb Q_N}{\mathrm d\mathbb P}
 \leq \exp(k_N),\qquad k_N=o(N).
\end{equation}
The independent grid displacement introduced below has the same conditional
law under both measures.

\begin{definition}[Projected address forest]\label{arm:address-forest}
Fix an integer $R\geq8$ and write $\lambda=\log R$.  Let $I$ be a
fixed interval of length $H$, where $H$ is a fixed constant multiple of the
base time unit $n$.  Let
\[
 J=\left\lfloor\frac{bN}{\lambda}\right\rfloor,
 \qquad \delta_j=H R^{-j},\quad 0\leq j\leq J.
\]
All retained time units are assumed to diverge.  Select a finite list of
markers $(t,\ell)$ with $t\in I$ and $0\leq\ell\leq J$.  The list,
including the exact positions, terminal levels, and all incidences, may
be a function of the entire word.  \emph{After this selection}, sample
$U$ uniformly on $[0,H)$, independently of everything already selected.
Use the nested half-open grids
\[
 \mathcal G_j(U)=
 \big\{[U+k\delta_j,U+(k+1)\delta_j):k\in\mathbb Z\big\}.
\]
A marker $(t,\ell)$ occupies the containing cell at each level $j\leq\ell$.
Let $V_j$ be the set of occupied cells, and put
\[
 m_j=|V_j|,\qquad T=\sum_{j=0}^Jm_j,
 \qquad A=\lambda T.
\]
The occupied cells, with the containing-cell parent relation, form the
\emph{projected address forest}.  Write $M$ for the total number of marker
entries, with multiplicity; entries terminating at an internal forest
vertex are retained as terminal labels there.

There are $h_0$ distinguished bases, each marked through level $J$.
At a probe level $j_p$, a witness list consists of finitely many chains
from these bases.  Every chain endpoint and intermediate position is
marked through $j_p$.  Every link is either a chord of one of the two
colors, or a time jump satisfying
\[
 |s-t|\leq K\delta_{j_p}.
\]
Here $K$ is a fixed finite list bound, enlarged to bound the number of
marker and edge entries per probe as well.  Thus each chain projects in
its entirety to every coarser level.  Exact endpoints within their
terminal cells will not be part of the address code.
\end{definition}

Only a bounded number of root cells can meet $I$; with the convention in
Definition~\ref{arm:address-forest}, there are at most two.  Their integer
names are determined by $I$ and $U$, so no unspecified macroscopic time
coordinate must be encoded.  All the sampling slots below are contained
in a fixed deterministic enlargement of $I$.  The bilateral word supplies
their chronological increments, including increments outside $I$ if a
slot meets an endpoint of that interval.

\begin{lemma}[Projected graph ranks]\label{arm:projected-rank}
At level $j$, retain all occupied boxes, and place the projected colored
chords between them, including isolated vertices.  Let $c_j$ be the
number of components of the union graph, let $b_j$ be the number of
occupied base boxes, and let $D_j$ count active jump entries whose two
endpoints lie in different level-$j$ boxes.  If $e_j$ is the sum of the
two monochrome graph ranks, then
\begin{equation}\label{arm:rank-component-bound}
 c_j\leq b_j+D_j,
 \qquad e_j\geq m_j-c_j\geq m_j-b_j-D_j.
\end{equation}
These counts are made before any sampling slots are discarded.
\end{lemma}
\begin{proof}
Include the projected jump edges temporarily.  Every occupied vertex
then has a chain to a base: a marker arises in a fully retained chain,
or is itself a base.  The enlarged graph therefore has at most $b_j$
components.  Removing the $D_j$ jump edges increases the component count
by at most $D_j$.  For the rank inequality, take a spanning forest of the
remaining union graph.  Its edges of either color form a forest in that
monochrome graph, so their number is bounded by that color's rank.
Summing the two bounds gives $m_j-c_j\leq e_j$.
\end{proof}

\begin{lemma}[Regular slots]\label{arm:regular-slots}
For every $\varepsilon>0$, one can choose a fixed sufficiently large
$R$ such that, conditional on any selected marker system, with
probability at least $7/8$ over $U$, all but
\begin{equation}\label{arm:regular-deletion}
 C_R M+\varepsilon T
\end{equation}
of its occupied vertex occurrences admit the following disjoint paired
slots.  At a retained vertex $v\in V_j$, let $t_v$ be the midpoint of
its cell and put $r_j=R^2\delta_j$.  The slots are
\begin{equation}\label{arm:paired-slots}
 S_v^-=[t_v-r_j,t_v-2r_j/R],
 \qquad
 S_v^+=[t_v+2r_j/R,t_v+r_j].
\end{equation}
All these intervals, for every retained vertex, level, and sign, are
pairwise disjoint.  Their locations are determined by the forest and
$U$.  The constant $C_R$ is independent of the positions, the number of
probe levels, and the finite exploration cutoff except through $M$.
\end{lemma}
\begin{proof}
Call a vertex critical if it is a root, has a number of children different
from one, or carries a terminal entry.  There are at most $M$ leaves and
at most $M$ terminating vertices.  For a finite forest,
\[
 \sum_{v:\,\deg^+(v)\geq2}(\deg^+(v)-1)
 =\#\{\text{leaves}\}-\#\{\text{roots}\},
\]
so there are at most $M$ branching vertices as well.  Thus at most $4M$
vertices are critical.  Fix an integer $g=4$ and discard all vertices at
forest distance at most $g$ from a critical vertex.  Since the degree is
at most $R+1$, this discards at most
\[
 C_RM,\qquad
 C_R=4\sum_{k=0}^g(R+1)^k.
\]

For a remaining vertex $v\in V_j$, its ancestor $Q_v$ at depth $j-g$
exists.  The intervening path is unary.  In particular, $v$ is the only
occupied depth-$j$ descendant of $Q_v$.  Additionally discard $v$ if
its cell is at distance less than $3r_j$ from $\partial Q_v$.
Call this a boundary deletion.  We first prove disjointness for the
vertices left by these two deletion rules.

Take distinct retained vertices $v\in V_j$ and $w\in V_k$, with
$k\geq j$.  If $k=j$, the cell of $w$ is outside $Q_v$, since $v$ is
its sole occupied depth-$j$ descendant.  Hence
$|t_w-t_v|\geq3r_j$, whereas the sum of outer slot radii is $2r_j$.
If $k>j$ and $w$ descends from $v$, then
\[
 |t_w-t_v|\leq\delta_j,\qquad r_k\leq r_j/R,
 \qquad
 \delta_j+r_j/R
 =r_j(R^{-2}+R^{-1})<2r_j/R.
\]
Both slots at $w$ therefore lie strictly inside the inner gap of the
slots at $v$.  If $k>j$ and $w$ does not descend from $v$, its occupied
ancestor at depth $j$ lies outside $Q_v$.  Thus $t_w$ is outside $Q_v$,
$|t_w-t_v|\geq3r_j$, and the sum of outer radii is at most
$r_j+r_j/R<2r_j$.  This proves disjointness in every case.  The two
slots at a single vertex are disjoint by their definition.

It remains to count boundary deletions without imposing any distribution
on the marker positions.  Condition on all these positions and terminal
levels.  At each depth $j$, let $s_j$ be the occupancy count for a fixed
auxiliary grid of width $\delta_j$, with no random displacement, and
let $S=\sum_j s_j$.  Each cell of either of two grids of the same width
meets at most two cells of the other grid.  It follows, for every $U$,
that
\begin{equation}\label{arm:grid-comparability}
 S/2\leq T(U)\leq2S.
\end{equation}
Choose one active marker in each occupied auxiliary cell.  If an
occupied primary cell at depth $j\geq g$ is boundary-deleted, a chosen
representative in an auxiliary cell meeting it lies within
$3r_j+3\delta_j$ of a depth-$(j-g)$ boundary.  The displacement modulo
$\delta_{j-g}$ is uniform.  The probability of this event for a fixed
representative is at most
\[
 \frac{2(3r_j+3\delta_j)}{\delta_{j-g}}
 \leq12R^{2-g}.
\]
An auxiliary cell meets at most two primary cells.  If $B(U)$ is the
number of boundary deletions, linearity of expectation gives
\[
 \mathbb E_U B(U)\leq24R^{2-g}S.
\]
By Equation~\eqref{arm:grid-comparability} and Markov's inequality,
\begin{equation}\label{arm:boundary-probability}
 \mathbb P_U\{B(U)>\varepsilon T(U)\}
 \leq\frac{48R^{2-g}}{\varepsilon}.
\end{equation}
For $g=4$, choose $R$ large enough that this is at most $1/8$.
Together with the structural deletions this proves
Equation~\eqref{arm:regular-deletion}.
\end{proof}

Figure~\ref{fig:regular-slots} illustrates the nested-slot separation and
the padding inside the unary ancestor.
\begin{figure}[tbp]
\centering
\begin{tikzpicture}[
  x=1cm,y=1cm,
  font=\small,
  >=Latex,
  coarse/.style={draw=blue!65!black,fill=blue!13,line width=.8pt},
  finer/.style={draw=teal!70!black,fill=teal!20,line width=.8pt},
  guide/.style={draw=black!35,densely dashed,line width=.45pt},
  dimension/.style={<->,draw=black!65,line width=.5pt},
  every node/.style={inner sep=2pt}
]
\node[anchor=west,font=\footnotesize] at (.35,4.18)
  {Cell nesting (magnified, schematic)};
\draw[draw=black!55,rounded corners=1.5pt] (.35,2.88) rectangle (14.05,3.88);
\node[anchor=west] at (.55,3.65) {$Q_v$};
\node[anchor=east,font=\footnotesize,text=black!65] at (13.85,3.65)
  {unary ancestor at depth $j-4$};
\draw[coarse] (5.85,3.04) rectangle (8.55,3.69);
\node[anchor=west,text=blue!65!black] at (6.0,3.37) {$v$};
\draw[finer] (7.26,3.14) rectangle (7.94,3.59);
\node[text=teal!70!black] at (7.60,3.37) {$w$};
\draw[dimension] (.55,3.13) -- (5.65,3.13)
  node[midway,above,font=\footnotesize] {$\ge 3r_j$};
\draw[dimension] (8.75,3.13) -- (13.85,3.13)
  node[midway,above,font=\footnotesize] {$\ge 3r_j$};

\fill[black!4] (5.75,-.70) rectangle (8.65,2.20);
\draw[guide] (5.75,-.70) -- (5.75,2.20);
\draw[guide] (8.65,-.70) -- (8.65,2.20);
\draw[->,draw=black!60] (.35,-.47) -- (14.05,-.47);
\node[anchor=east,font=\footnotesize] at (14.05,-.76) {time};

\draw[coarse] (1.4,1.03) rectangle (5.75,1.45);
\draw[coarse] (8.65,1.03) rectangle (13.0,1.45);
\node[text=blue!65!black] at (3.575,1.24) {$S_v^-$};
\node[text=blue!65!black] at (10.825,1.24) {$S_v^+$};
\node[anchor=south] at (1.4,1.62) {$t_v-r_j$};
\node[anchor=south] at (5.75,1.62) {$t_v-2r_j/R$};
\node[anchor=south] at (8.65,1.62) {$t_v+2r_j/R$};
\node[anchor=south] at (13.0,1.62) {$t_v+r_j$};
\draw[draw=blue!65!black,line width=.65pt] (7.20,.85) -- (7.20,1.53);
\fill[blue!65!black] (7.20,1.24) circle (1.6pt);
\node[anchor=south,text=blue!65!black] at (7.20,1.60) {$t_v$};

\draw[finer] (6.505,.05) rectangle (7.04875,.47);
\draw[finer] (7.41125,.05) rectangle (7.955,.47);
\node[anchor=south,text=teal!70!black] at (6.77,.52) {$S_w^-$};
\node[anchor=south,text=teal!70!black] at (7.69,.52) {$S_w^+$};
\draw[draw=teal!70!black,line width=.65pt] (7.23,-.14) -- (7.23,.53);
\fill[teal!70!black] (7.23,.26) circle (1.6pt);
\draw[draw=teal!70!black,line width=.45pt] (7.23,-.14) -- (7.67,-.30);
\node[anchor=west,text=teal!70!black] at (7.67,-.30) {$t_w$};
\node[anchor=west,font=\footnotesize] at (.55,.26)
  {finer descendant, $k>j$};
\draw[dimension] (5.75,-1.02) -- (8.65,-1.02);
\node[anchor=north,font=\footnotesize] at (7.20,-1.08)
  {coarse central gap};
\end{tikzpicture}
\caption[Disjoint regular sampling slots]{%
Regular sampling slots from Lemma~\ref{arm:regular-slots}.
The inset shows cell containment; the two slot rows share one time axis
and are vertically separated for visibility.  Here $r_j=R^2\delta_j$,
the unary ancestor $Q_v$ lies at depth $j-4$, $v$ is its sole occupied
depth-$j$ descendant, and $\operatorname{dist}(v,\partial Q_v)\ge3r_j$.
For a finer descendant $w\subset v$ at depth $k>j$,
$|t_w-t_v|+r_k\le\delta_j+r_j/R
=r_j(R^{-2}+R^{-1})<2r_j/R$;
thus both finer slots lie strictly inside the coarse central gap,
although they need not lie inside the cell $v$.
For an incomparable retained node at depth $k\ge j$,
$|t_w-t_v|\ge3r_j>r_j+r_k$.
These inequalities give disjointness in both cases.}
\label{fig:regular-slots}
\end{figure}


\begin{remark}[Other fixed test bands]\label{arm:other-test-bands}
The same construction applies to a fixed finite family of test bands
contained in $a_*r\leq |t-t_v|\leq A_*r$, where
$0<a_*<A_*<\infty$.  First refine the location boxes until their width
is at most $a_*r/4$, and sample generations far enough apart that the
next outer radius is at most $a_*r/4$.  This uses a fixed larger grid
ratio or a fixed residue class of generations; an independent uniform
choice of the residue preserves uniform sampling of levels.  Increase
the boundary padding beyond $2A_*r$ and then increase the fixed ancestor
depth $g$ if necessary.  A finer descendant's bands are now inside the
coarser inner gap, while an incomparable descendant lies outside the
padded ancestor.  The preceding disjointness proof applies verbatim to
these two explicit inequalities.  The structural and boundary counts
have the same form, with constants depending on the fixed bands.
For integer word times, rounding slot boundaries inward preserves
disjointness and changes normalized lengths by a vanishing amount.
\end{remark}

\begin{lemma}[Integrated connectivity errors]\label{arm:jump-errors}
For each fixed $K,R$, a jump entry at terminal level $\ell$ contributes
at most $C_{K,R}<\infty$ to the expected sum of its separated coarse
projections.  Consequently, if $M_{\mathrm{jump}}$ is the number of jump
entries,
\begin{equation}\label{arm:jump-sum}
 \mathbb E_U\sum_{j=0}^JD_j
 \leq C_{K,R}M_{\mathrm{jump}}.
\end{equation}
If the bases lie in an interval of length at most $C H e^{-N}$ and
$j_N=\lfloor N/\lambda\rfloor$, then also
\begin{equation}\label{arm:base-merger-error}
 \mathbb E_U\sum_{j=0}^{j_N}(b_j-1)\leq C'_{C,R}.
\end{equation}
For a sparse family with $P=O(N/L)$ probes, $L\longrightarrow\infty$,
and fixed list bound $K$, there is a deterministic $d_N=o(N)$ such that,
uniformly in the selected witness data,
\begin{equation}\label{arm:good-connectivity}
 \mathbb P_U\left\{
 \sum_{j=0}^JD_j+
 \sum_{j=0}^{j_N}(b_j-1)>d_N\right\}=o(1).
\end{equation}
Moreover $C_RM=o(N)$.  Together with
Lemma~\ref{arm:regular-slots}, these assertions hold simultaneously on
a set of displacements of conditional probability at least $1/2$ for
all sufficiently large $N$.
\end{lemma}
\begin{proof}
Two fixed times at distance $d$ lie in different uniformly displaced
cells of width $\delta_j$ with probability at most
$\min\{1,d/\delta_j\}$.  Therefore a jump of length at most
$K\delta_\ell$ has expected total separated projections at most
\[
 \sum_{j=0}^\ell\min\{1,K R^{j-\ell}\}
 \leq \sum_{a=0}^\infty\min\{1,K R^{-a}\}
 =:C_{K,R}<\infty.
\]
Linearity proves Equation~\eqref{arm:jump-sum}.  For
Equation~\eqref{arm:base-merger-error}, list the bases in increasing
order and join consecutive bases by artificial edges for counting only.
At every depth, $b_j-1$ is bounded by the number of these edges that
cross a cell boundary.  Their total Euclidean length is at most
$CHe^{-N}$.  Hence
\[
 \mathbb E_U(b_j-1)\leq C e^{-N}R^j,
 \qquad
 \sum_{j=0}^{j_N}C e^{-N}R^j\leq \frac{CR}{R-1}.
\]
For fixed $K$, both $M$ and $M_{\mathrm{jump}}$ are bounded by a fixed
multiple of $P+h_0$.  Choose $d_N=o(N)$ so slowly that
$(P+h_0+1)/d_N\longrightarrow0$.  Markov's inequality gives
Equation~\eqref{arm:good-connectivity}; the bound $C_RM=o(N)$ follows
from $P=o(N)$.  The last assertion follows by intersecting this event
with the event of Lemma~\ref{arm:regular-slots}.
\end{proof}

\begin{lemma}[Stability under an independent grid displacement]
\label{arm:grid-partition-stability}
Fix a marker system with $M$ entries before sampling its displacement.
At level $j$, partition its active entries without using a grid: first
group entries at the same time, then join consecutive distinct active
positions whenever their gap is strictly less than $\delta_j$.
Denote the resulting partition by $\mathcal P_j^*$, and let
$\mathcal P_j(U)$ be the partition into occupied cells of the displaced
grid.  Put
\[
 \Delta(U)=\#\{0\le j\le J:\mathcal P_j(U)\ne\mathcal P_j^*\}.
\]
For an independent uniform displacement,
\begin{equation}\label{arm:grid-partition-error}
 \mathbb E_U[\Delta(U)\mid\text{marker system}]
 \le \frac{2R}{R-1}M.
\end{equation}
At every level where the partitions agree, the entire projected
colored graph, including its base labels, is the same up to cell names.
In particular, if $M\le M_N=o(N)$, two independent uniform
displacements give the same projected graphs outside $o(N)$ levels
with probability tending to one, uniformly over the selected marker
system.
\end{lemma}

\begin{proof}
If $M=0$, all partitions are empty and the assertions are immediate.
Keep the labels and incidences of coincident entries when grouping
their positions.  As $j$ increases, entries and hence distinct active
positions can only disappear.  Initially there are at most $M-1$
consecutive-position pairs; deleting one position creates at most one
new pair.  Thus at most $2M-1$ distinct consecutive pairs occur over
all levels.  Deleting all positions with a common terminal level one
at a time only enlarges this bound.

Every grid partition refines $\mathcal P_j^*$: positions in the same
half-open cell have total span strictly less than $\delta_j$, and so
do all their successive gaps.  The partitions can differ only if a
consecutive active pair with gap $d<\delta_j$ is split by a cell
boundary.  For a fixed pair this has probability $d/\delta_j$, since
$U$ modulo $\delta_j$ is uniform.  For each positive gap,
\[
 \sum_{j:\,\delta_j>d}\frac d{\delta_j}
 \le \frac R{R-1}.
\]
The same bound holds after restricting to the levels at which the
pair is consecutive and active.  A union bound over pairs and levels
proves Equation~\eqref{arm:grid-partition-error}.

All active edge incidences and base labels were fixed before the grid,
so equal marker partitions give identical projected graph data.
Choose a deterministic $q_N=o(N)$ with $M_N/q_N\to0$.
Markov's inequality bounds the conditional probability of
$\Delta(U)>q_N$ by $2RM_N/((R-1)q_N)=o(1)$.  Apply this to each of
two independent displacements.  Outside that exceptional event they
both agree with $\mathcal P_j^*$ except at at most $2q_N$ levels.
The estimate is conditional on arbitrary marker data, so it also
holds under any selection law sampled before the displacements.
\end{proof}

\begin{remark}[Why the order of selection matters]
The witnesses in Definition~\ref{arm:address-forest} must precede $U$.
For example, take two terminal positions $U-\eta$ and $U+\eta$ with
$\eta\ll\delta_J$.  They occupy distinct unary branches on opposite
sides of the persistent boundary $U$ at all $J+1$ levels, but the
same-level paired slots overlap.  At least half the $2(J+1)$ nodes
must be deleted, although there are only two terminal entries.  This
contradicts a bound $C_RM+\varepsilon T$ with $\varepsilon<1/2$ for
large $J$.  Our conditional estimates avoid this obstruction.  We keep
$U$ independent in the probability argument, rather than treating a
sample-dependent favorable displacement as free fixed information.
\end{remark}

\Needspace{8\baselineskip}
\subsection{Counting the address codes}

\begin{lemma}[Address code count]\label{arm:address-codes}
Fix $U$ and the distinct probe depths $j_1,\ldots,j_P$.  Include the
base entries in the probe list at depth $J$, adding that depth if needed,
and enlarge $K$ to cover those entries.  The resulting probe depths
remain distinct.  Suppose each list has at most $K$ marker and edge entries.  There is a coding of the
projected forests and their list incidences with the following bound,
uniformly in the total occupancy $T$:
\begin{align}
 \log\#\{\text{codes of occupancy }T\}
 &\leq(\lambda+C_0)T
   +KP\log(1+T/P)
   +CPK\log(K+1)+O(1).\label{arm:address-code-count}
\end{align}
The constant $C_0$ does not depend on $K$.  A family of additional
annotations with at most $C_a^{m_j}$ choices at depth $j$, for a fixed
$C_a$, only increases $C_0$ by $\log C_a$.
In particular, for fixed $K$ and $P=o(N)$,
\begin{equation}\label{arm:address-code-leading}
 \log\#\{\text{codes of occupancy }T\}
 \leq(\lambda+C_0)T+o(T+N),
\end{equation}
uniformly in $T$.  A sufficient condition for a growing bound $K=K_N$
and probe stride $L$ is
\begin{equation}\label{arm:growing-list-condition}
 \frac{K_N\log(K_NL)}{L}\longrightarrow0.
\end{equation}
The codes can also be assigned subprobability weights at an additional
cost $2\log(T+1)+O(1)$.
\end{lemma}
\begin{proof}
An ordered forest with $T$ vertices and at most two roots has at most
$C4^T$ possible shapes.  One way to see this is to attach a new root
above its roots and use the depth-first parenthesis word of the
resulting ordered tree; its length is $2(T+1)$, and a balanced
parenthesis word determines the tree.  Each nonroot cell has one of
$R$ child offsets in its parent, so the shape has at most $R^T$
geometric assignments.  The occupied root cells have a bounded presence
mask among the cells meeting $I$, whose names are already fixed by
$I,U$.  This proves the bound $C(4R)^T$ for the forest before labels.
Its vertices' depths follow from its parent relation.

At probe depth $j_p$, name a position by its ordered index among the
$m_{j_p}$ occupied cells; use one extra index for an absent list entry.
Thus the position entries have at most $(1+m_{j_p})^K$ possibilities.
Repeated indices record coincident positions.  Names in different
probe lists do not require global forest indices: each local index
specifies a forest vertex, whose ancestors are already determined by
the forest.  A correspondence to a distinguished base uses that base's
fixed name, and its depth-$J$ address determines its coarser addresses.
The list's chain order, colors, endpoint incidences, and base names have
at most $\exp(CK\log(K+1))$ possibilities.  The count includes absent
entries and shorter lists.  It never specifies an exact endpoint
within its terminal cell.

Multiply these choices and use the distinctness of probe depths to
obtain $\sum_p m_{j_p}\leq T$.  Concavity gives
\[
 \sum_{p=1}^P\log(1+m_{j_p})
 \leq P\log(1+T/P),
\]
which proves Equation~\eqref{arm:address-code-count}.  The same estimate
bounds the maximum number of list decorations over all forest shapes,
so variation of the individual $m_j$ between shapes causes no extra
factor.

For $0<\eta\leq K$ and $z\geq0$,
\[
 K\log(1+z)\leq\eta z+K\log(K/\eta).
\]
Indeed the maximum of $\log(1+z)-(\eta/K)z$ occurs at
$z=K/\eta-1$ and is at most $\log(K/\eta)$.  Hence
\begin{equation}\label{arm:address-code-formbound}
 \log\#\{\text{codes of occupancy }T\}
 \leq(\lambda+C_0+\eta)T
   +KP\log(K/\eta)+CPK\log(K+1)+O(1).
\end{equation}
For fixed $K$, let $\eta$ tend to zero sufficiently slowly while
$P=o(N)$.  This yields Equation~\eqref{arm:address-code-leading}
uniformly in $T$.  Under Equation~\eqref{arm:growing-list-condition},
for example $\eta_N=1/\log L$ for large $L$ makes both remainders in
Equation~\eqref{arm:address-code-formbound} equal to $o(N)$.
No condition $L\gg\log N$ is used.  If the probe depths themselves
must be encoded among the $J+1$ levels, their extra cost is at most
$\log\binom{J+1}{P}\leq P\log(e(J+1)/P)=o(N)$ for
$P=O(N/L)$; alternatively their independent mesh displacement can be
conditioned on.

Finally give $T$ a weight proportional to $(T+1)^{-2}$ and distribute
that weight uniformly among its finite enumeration of codes.  The
weights sum to at most one, and the negative logarithm of each weight
is the code count bound plus $2\log(T+1)+O(1)$.
\end{proof}

We call a displacement \emph{favorable} if it satisfies the boundary
bound of Lemma~\ref{arm:regular-slots} and the connectivity bound of
Lemma~\ref{arm:jump-errors}.  For fixed $R,K$, its conditional
probability given the selected witness system is at least $1/2$ for
large $N$.  At a favorable displacement the structural deletions are
$o(N)$ and
\begin{equation}\label{arm:total-rank-lower}
 E:=\sum_{j=0}^Je_j
 \geq T-h_0(J+1)-d_N.
\end{equation}
For earlier scales, Equation~\eqref{arm:base-merger-error} also gives
the stronger replacement of $h_0$ by one, up to an integrated $o(N)$
error.  The simpler bound in Equation~\eqref{arm:total-rank-lower}
suffices for the following occupancy estimate.


The code now specifies every box, component, and sampling slot.  It remains
to bound the probability that its required contacts occur.  We first
control the sum of contour oscillations across its prescribed intervals.

\subsection{Oscillations and the probability of a fixed code}

\begin{lemma}[An exponential sum bound for prescribed intervals]
\label{arm:oscillation-sum}
Consider a deterministic finite family $\mathcal I$ of intervals whose
lengths belong to a geometric family with ratio $R>1$.  Suppose that
at each length at most $K$ intervals cover a point.  There are constants
$c,C>0$, depending on $R,K$ and the raw-word law, such that
\[
 \mathbb P\left\{
  \sum_{I\in\mathcal I}\frac{\operatorname{osc}_I Z}{\sqrt{|I|}}
       >u|\mathcal I|\right\}
 \le \exp\{-cu|\mathcal I|+C|\mathcal I|\}.
\]
Here coordinate oscillations can be summed, and the estimate holds for
all exterior typings.  Integer intervals of bounded length can be
handled by increasing the constant.  If $K$ is fixed and $R\ge2$, the
constants can be chosen uniformly in $R$.
\end{lemma}

\begin{proof}
Cut the union of the intervals at all their endpoints.  There are at
most $2|\mathcal I|$ nonempty atoms $J$.  Proposition~\ref{arm:raw-exponential}
gives independent dominating variables
$Y_J=2M_{|J|}/\sqrt{|J|}$, with a common exponential-moment bound, for
the oscillations on these disjoint fresh raw pieces.  Subadditivity
of oscillation gives
\[
 \sum_{I\in\mathcal I}\frac{\operatorname{osc}_I Z}{\sqrt{|I|}}
 \le\sum_J a_JY_J,
 \qquad
 a_J=\sqrt{|J|}\sum_{I\supset J}|I|^{-1/2}.
\]
If $I_0$ is the shortest interval containing $J$, then
\[
 a_J\le K\sqrt{|J|/|I_0|}
             \sum_{v\ge0}R^{-v/2}
 \le \frac K{1-R^{-1/2}}.
\]
Apply the exponential-moment bound with a parameter divided by this
last constant, and use Markov's inequality.  This proves the claim.
The bound on $a_J$ is uniform for $R\ge2$, proving the asserted
uniformity.
\end{proof}

\begin{proposition}[Fixed-code arm test]\label{arm:fixed-code-test}
Let $\lambda=\log R$.  In the finite shifted address model, fix a code,
let $m_j$ be its number of occupied boxes at level $j$, and let $e_j$
be the sum of its two monochrome ranks.  Put
\[
 T=\sum_jm_j,\qquad E=\sum_je_j,
\]
and let $M$ be the number of terminal entries.  For each $\zeta>0$, one
can choose a fixed sufficiently large $R$ and sufficiently small fixed
deletion tolerances so that, as the minimum sampling radius tends to
infinity,
\begin{equation}\label{arm:fixed-code-probability}
 \mathbb P\{\hbox{the code's chord constraints are realizable}\}
 \le \exp\{-\beta\lambda E+\zeta\lambda T+C_{R,\zeta}M\}.
\end{equation}
The statement is for codes with the regular-slot deletion bound of
Lemma~\ref{arm:regular-slots}.  It is uniform over such finite codes
and over arbitrary exact chord endpoints inside their recorded boxes.
Small fixed errors in sampling radii or inward integer rounding are
absorbed by increasing the minimum permitted radius.
\end{proposition}

\begin{proof}
Choose a deterministic time before all the code's representative
points and sampling slots, and anchor both contour coordinates at
zero there.  A common coordinate translation changes none of the
chord constraints or slot tests.  With this convention the heights
observed while exposing the finite word window are measurable in
its chronological raw filtration, even when the original bilateral
contours had been anchored at time zero inside that window.

For each coordinate and each projected chord component, apply the
sorted comparison lemma to the occupied intervals of that component.
It gives a comparison between consecutive occupied intervals, one per
unit of monochrome rank, with a total height-discrepancy bound of at
most six times the sum of the relevant box oscillations.  A box is
used in at most one component of each color.  Therefore the sum of
all normalized comparison discrepancies is controlled by a fixed
multiple of the normalized box oscillations.

Include also the intervals between each regular-node representative midpoint and
its adjacent sampling-slot endpoint.  Their lengths are comparable
to $r/R$ at a node of sampling radius $r$.  The same-scale overlaps
are bounded independently of $R$: regular centers are separated at
the scale $r$, whereas the box intervals themselves have disjoint
interiors.  Thus Lemma~\ref{arm:oscillation-sum} applies, separately
to a bounded number of geometric length families.  For fixed
$0<\epsilon<\delta$, outside an event of probability at most
\[
 \exp\{-cR^\epsilon T+CT\},
\]
the sum of all these normalized oscillations is at most $R^\epsilon T$.
Discard each comparison whose assigned discrepancy exceeds
$R^\delta\sqrt{r/R}$, and each node-side with an excessive local
oscillation.  At most $C R^{\epsilon-\delta}T$ comparison incidences
are lost.  The originally irregular nodes lose at most
$C\epsilon_0T+C_RM$ incidences: after sorting, a node has at most
one predecessor and one successor per color.  Choosing $R$ large
makes the first two relative losses as small as required.

All remaining required slot conditions now have fixed deterministic
tolerance $C R^\delta\sqrt{r/R}$.  On a left slot the target wall is
the already observed representative height in the previous occupied
box of the same component, minus that tolerance.  The comparison
bound makes survival above this wall and landing near it necessary.
The previous box lies strictly before the slot, by regular-node
separation.  On a right slot use the current midpoint height minus
the same tolerance; that height is already observed before this
slot.  Its initial discrepancy is small by the retained local
oscillation condition.  We do not choose a random wall using the
unobserved actual comparison error.

Expose all slots in chronological order.  Lemma~\ref{arm:regular-slots}
makes them disjoint, and the intervening raw words can also be
exposed in chronological order.  Each slot is fresh conditional on
this history.  If $k$ colors are incident on a retained left slot,
Lemma~\ref{arm:lattice-slots} gives logarithmic charge at least
\[
 (k/2+d_k-O(\delta)-o(1))\lambda.
\]
For a right slot it gives
$(d_k-O(\delta)-o(1))\lambda$.  The $o(1)$ is uniform once the
minimum radius is sufficiently large for the fixed $R$.
Summed incoming incidences give $E/2$ before deletions.  There are
$2E$ side incidences in total, and
$d_k\ge k d_2/2$ for $k=0,1,2$.  The undeleted leading survival
charge is therefore at least $d_2E$.  Deletions lose only a bounded
charge per incidence.  The choices of deleted incidences have at
most $\exp(CT)$ possibilities, which are summed rather than
conditioned upon.

Choose $\delta$ sufficiently small in terms of $\zeta$, then choose
$R$ large enough that $C/\lambda$, the deletion fractions, and the
fixed Brownian constants divided by $\lambda$ fit within the
remaining error budget.  The oscillation-exception exponent can
simultaneously be made larger than $3\beta\lambda T$ in absolute
value, since $R^\epsilon/\log R\to\infty$.  As $E\le2T$, that
exception is absorbed into Equation~\eqref{arm:fixed-code-probability}
after an arbitrarily small enlargement of its error budget.  Finally
take the minimum radius large enough for the uniform lattice-slot
estimates.  Chronological multiplication and the finite sum over
discard patterns give the proposition.
\end{proof}


\subsection{The address occupancy bound}

The fixed-code probability bound outweighs the address count because
$\beta>1$.  We combine them before returning to the independent grid,
so a favorable displacement is never treated as uncharged selected data.

\begin{proposition}[Exponential occupancy bound]\label{arm:occupancy-tail}
Put $\beta=1/2+d_2>1$.  Assume the fixed-code test of
Proposition~\ref{arm:fixed-code-test} in the following quantitative form:
for every sufficiently small $\zeta>0$, choose a fixed boundary
tolerance sufficiently small and a fixed $R$ sufficiently large so
that, once the minimum radius is sufficiently large, every
\emph{favorable admissible} decorated code $\mathfrak c$ satisfies
\begin{equation}\label{arm:required-fixed-code-bound}
 \mathbb P\{\mathfrak c\text{ is realized}\}
 \leq\exp\bigl\{-\beta\lambda E(\mathfrak c)
          +\zeta\lambda T(\mathfrak c)
          +C_{R,\zeta}M(\mathfrak c)\bigr\}.
\end{equation}
This estimate is required uniformly in the fixed displacement $U$.
The realization event includes all the actual marked chord constraints.
The fractional boundary deletions are absorbed into the $\zeta\lambda T$
term, and the structural deletions into the explicit terminal term,
using the bounded number of sorted-chain charges at each deleted
vertex.  The fixed-code estimate includes its oscillation exceptions.
No such estimate is asserted here for codes whose slots are
predominantly irregular.

Then there are $C,c,C_1>0$, independent of any fixed finite exploration
cutoff, such that for all sufficiently large $N$ and all $u\geq0$,
\begin{equation}\label{arm:occupancy-exponential-tail}
 \mathbb P\{A\geq CN+u\}
 \leq C_1\exp(-cN-cu).
\end{equation}
The estimate holds for any word-dependent selection of witness lists
made before $U$.  It also holds for a fixed reference grid, and with
$A$ replaced by its maximum over grid displacements.  The threshold
in ``sufficiently large'' may depend on the fixed exploration cutoff.
If Equation~\eqref{arm:address-tilt} holds, then
\begin{equation}\label{arm:occupancy-tilted-mean}
 \mathbb E_{\mathbb Q_N}A=O(N).
\end{equation}
The corresponding code has length at most a deterministic constant
times $N$ with $\mathbb Q_N$-probability tending to one.  At an
independently and uniformly sampled level $j\in\{0,\ldots,J\}$,
$\mathbb E_{\mathbb Q_N}m_j=O(1)$.
\end{proposition}
\begin{proof}
Fix $U$.  Use Equation~\eqref{arm:address-code-formbound} with a small
fixed $\eta>0$ and then let the sparse terminal error be $o(N)$.
For every favorable code of total occupancy $T$, combine
Equations~\eqref{arm:required-fixed-code-bound} and
\eqref{arm:total-rank-lower}.  Summing over its possible codes gives
\begin{align*}
 &\mathbb P\{\text{some favorable realized code has occupancy }T\}\\
 &\quad\leq
 \exp\bigl\{
 -[(\beta-1-\zeta)\lambda-C_0-\eta]T
 +\beta\lambda h_0(J+1)
 +\beta\lambda d_N+C_{R,\zeta}M+o(N)
 \bigr\}.
\end{align*}
Choose fixed $R$ large enough and then the error tolerances small enough
that
\[
 \delta:=\beta-1-\zeta-(C_0+\eta)/\lambda>0.
\]
Since $\lambda(J+1)=bN+O(1)$, $d_N=o(N)$, and $M=o(N)$, there is a fixed
$B<\infty$ such that the last probability is at most
\begin{equation}\label{arm:fixed-occupancy-union}
 \exp\{-\delta\lambda T+BN\}.
\end{equation}
Here $B$ and $\delta$ are independent of the fixed list cutoff: its
contribution was confined to the sparse $o(N)$ terms.  Geometric summation
over $\lambda T\geq x$ shows, uniformly in $U$,
\[
 \mathbb P\{\text{favorable and }A\geq x\}
 \leq C_2\exp(BN-\delta x).
\]
For $C_*=(B+2)/\delta$, this gives
\begin{equation}\label{arm:favorable-tail}
 \mathbb P\{\text{favorable and }A\geq C_*N+u\}
 \leq C_2\exp(-2N-\delta u).
\end{equation}
The same bound holds after integrating over the independent $U$.

The event of a favorable displacement has only a positive conditional
probability; we do not give its complement an exponential estimate.
Instead compare all shifts of the same realized witness system.  The
same-width cell intersection argument gives
\[
 T(U)\leq2T(V)\qquad\text{for all displacements }U,V.
\]
Let $A_{\max}=\max_U\lambda T(U)$.  This maximum is measurable:
there are finitely many markers and levels, so only finitely many
marker--boundary crossings can change the occupancy as $U$ varies in
$[0,H)$.  If $A_{\max}\geq x$, every displacement has $A\geq x/2$,
and at least half of the displacements are favorable.  Conditional on
the complete witness system, this proves
\[
 \mathbf 1_{\{A_{\max}\geq x\}}
 \leq2\mathbb P_U\{\text{favorable and }A\geq x/2
                      \mid\text{witness system}\}.
\]
Integrate and apply Equation~\eqref{arm:favorable-tail}.  Increasing
the constant multiplying $N$ proves
Equation~\eqref{arm:occupancy-exponential-tail} for $A_{\max}$, and
therefore for every particular grid.  All word-dependent choices have
been included in the finite code enumeration.  The auxiliary favorable
shift has never been selected as unencoded sample-dependent information.

Under Equation~\eqref{arm:address-tilt}, the tail is bounded above by
\[
 \mathbb Q_N\{A\geq CN+u\}
 \leq\min\{1,C_1\exp(k_N-cN-cu)\}.
\]
Since $k_N=o(N)$, integration in $u$ gives
\[
 \mathbb E_{\mathbb Q_N}A
 \leq CN+C_1c^{-1}\exp(k_N-cN)=O(N).
\]
Equation~\eqref{arm:address-code-formbound}, including its optional
integer-size code, then gives the stated deterministic $O(N)$ code
cutoff outside an event of probability tending to zero.  Finally,
for an independently uniform level,
\[
 \mathbb E_{\mathbb Q_N}m_j
 =\frac{\mathbb E_{\mathbb Q_N}T}{J+1}=O(1),
\]
since $J+1$ is comparable to $N$ for fixed $R,b$.
\end{proof}

\begin{remark}[Finite cutoffs and the order of limits]
Fix the desired error tolerance first, then choose $R$ and the padding
constants.  For each fixed finite exploration bound, let $N$ and the
minimum time unit diverge, and let the probe stride satisfy
$L\longrightarrow\infty$ and $L=o(N)$.  Structural, jump, and terminal
code errors are then $o(N)$.  A subsequent slow diagonal may increase
the exploration and list bounds, retaining
Equation~\eqref{arm:growing-list-condition} and the vanishing jump
error.  For example, the constant in Equation~\eqref{arm:jump-sum}
is $O(1+\log_R K)$, so that the displayed growing-list condition also
controls that error when the jump tolerance is bounded by $K_N$.
The coefficient in the occupancy bound remains fixed; only the speed
of the diagonal changes.  The construction counts all projected
addresses, including intermediate markers from different probe levels.
It assumes no sparsity of the original candidate locations at a finest
scale.  Producing a saturated, jointly realizable finite-line system
from the occupancy bound is a separate step.
\end{remark}
\subsection{The credit contradiction and entropy of viewed blocks}

For a connected component $C$ of a union of two colored graphs, let
$e(C)$ be the sum of the ranks of the two colored graphs restricted to
$C$, where rank is the number of vertices minus the number of connected
components, including isolated vertices.  We say that $C$ has a
\emph{rank cycle} when $e(C)\ge |C|$, and that it is \emph{acyclic}
when $e(C)=|C|-1$.  The union spanning-tree argument in
Lemma~\ref{arm:projected-rank} shows that these alternatives exhaust the
possibilities.  Thus an acyclic component has exactly the minimum total
rank needed for connectivity.

\begin{proposition}[Rank credit]\label{arm:rank-credit}
Let $\varepsilon=\beta-1>0$.  Consider the projected witness systems
of the address construction on logarithmic levels $[0,bN]$, with
$h$ bases in an interval of length $O(ne^{-N})$.  Suppose their
pairwise time distances exceed the largest address-cell width used
on $[aN,bN]$, where $1<a<b$; thus they occupy distinct boxes there
for every displacement.  All selected
laws satisfy $\mathbb Q_N\le e^{o(N)}\mathbb P$.  Then the following
combination of events has probability tending to zero: outside a
vanishing fraction of levels in $[aN,bN]$, every connected
component containing a base has a rank cycle.

Here $a-1$ is chosen sufficiently small relative to
$\varepsilon(b-a)$, and then $h$ sufficiently large, with all choices
fixed before the scaling limit.  Witness-list cutoffs may increase
only along the sparse diagonals allowed by the address construction.
\end{proposition}

\begin{proof}
At a level with $m$ occupied boxes, $c$ union components, and total
monochrome rank $e$, one has
\[
 e\ge m-c,\qquad m-\beta e\le\beta c-\varepsilon m.
\]
The integrated component excess caused by projected jumps is $o(N)$
on favorable grids.  Before level $N$, all base cells can be charged
as one component up to $o(N)$; hence the possible positive credit
there is at most $C_0N+o(N)$, independently of $h$.  On $[N,aN]$,
there are at most $h$ base components, so positive credit is at most
$C_1(a-1)hN+o(N)$.  The constants include fixed endpoint-rounding
and grid-step changes.

At every level on which all base components have rank cycles,
$e\ge m$ after the negligible connectivity errors.  The separated
bases occupy at least $h$ boxes on $[aN,bN]$, so the credit there is
at most $-\varepsilon h$ per unit log length.  At exceptional levels
the positive credit can still be bounded by a constant times the
number of acyclic base components; a bound on the maximum $m$ is
not needed.  Proposition~\ref{arm:occupancy-tail} absorbs all small
relative errors in the integrated occupancy into an arbitrarily
small multiple of $N$.

Choose $a-1$ such that
$C_1(a-1)<\varepsilon(b-a)/4$, and then choose $h$ with
$C_0<\varepsilon(b-a)h/4$.  The total leading address budget minus
arm charge is then at most $-cN$ for some $c>0$, even after the
permitted errors.  Summing Equation~\eqref{arm:fixed-code-probability}
over the length-weighted codes gives an exponentially small
probability for the stated cycle event \emph{and a favorable
displacement}, under $\mathbb P$.  Multiplying by $e^{o(N)}$ gives
the same vanishing conclusion under $\mathbb Q_N$.

We now remove favorability without assuming that cyclicity is
invariant under a change of grid.  Let $\eta_N\downarrow0$ and denote
by $E_N(U)$ the cycle event with at most $\eta_NN$ exceptional
generations for the independent displacement $U$.  Generation count
and logarithmic length differ by the fixed factor $\lambda$.
Lemma~\ref{arm:grid-partition-stability} supplies $q_N=o(N)$ and a
uniform conditional error $\alpha_N=o(1)$ such that
$\Delta(U)\le q_N$ outside that error.  Let $E_N^*$ be the corresponding
cycle event for the grid-independent partitions $\mathcal P_j^*$,
with at most $\eta_NN+q_N$ exceptions.  Then
\[
 \mathbb Q_N(E_N(U))\le\mathbb Q_N(E_N^*)+\alpha_N.
\]
Sample a second independent displacement $V$.  Conditional on the
marker system, $V$ is favorable and has $\Delta(V)\le q_N$ with
probability at least $1/2-\alpha_N$.  On $E_N^*$ this gives the
projected cycle event $E_N'(V)$ with at most $\eta_NN+2q_N=o(N)$
exceptional generations.  Consequently
\[
 (1/2-\alpha_N)\mathbb Q_N(E_N^*)
 \le \mathbb Q_N(E_N'(V)\cap\{V\text{ favorable}\}).
\]
The preceding favorable-grid estimate applies to this right side.
Indeed, on each added exceptional level the positive credit is at
most $\beta(h+D_j)$, by Equation~\eqref{arm:rank-component-bound};
the extra levels and the favorable-grid sum $\sum_jD_j$ are both
$o(N)$.  No maximum-occupancy bound is needed for this error.  The
right side therefore tends to zero, proving the claim for $U$.
The conditional grid estimates are uniform in the preselected data;
their errors are not multiplied by the selected density.
\end{proof}

\begin{lemma}[Entropy allocation for a finite address code]
\label{arm:entropy-allocation}
Let $W$ have a product raw-word law $\mathbb P$, including an
independent past if needed.  Let $\mathbb Q\le e^\kappa\mathbb P$,
and let $C$ be a countably valued code chosen using $W$ and independent
auxiliary randomness.  Suppose $C$ has code lengths $\ell(c)$ with
$\sum_c e^{-\ell(c)}\le1$.  For each fixed code, let finitely many
disjoint deterministic word slots be ordered chronologically.  Then
the sum of their expected conditional relative entropies with respect
to fresh iid raw blocks, averaged over $C$, is at most
\begin{equation}\label{arm:entropy-total}
 \kappa+\mathbb E_{\mathbb Q}\ell(C).
\end{equation}
Auxiliary continuous grid shifts can be kept as independent side
information and integrated; no code length for such an independent
shift is charged.
\end{lemma}

\begin{proof}
Let $\pi_c=e^{-\ell(c)}$ and place any missing mass at an unused code.
Compare the joint law of $(C,W)$ under $\mathbb Q$ to the law in
which $C$ has probabilities $\pi_c$ and $W$ independently has law
$\mathbb P$.  The chain rule gives
\[
 D(\mathbb Q_{C,W}\Vert\pi\otimes\mathbb P)
 =D(\mathbb Q_W\Vert\mathbb P)
   +\mathbb E_{\mathbb Q_W}
        D(\mathbb Q_{C\mid W}\Vert\pi)
 \le\kappa+\mathbb E_{\mathbb Q}\ell(C),
\]
because the conditional Shannon entropy of $C$ is nonnegative.
Disintegrate instead first over $C$ and then over the chronological
raw pieces.  Product structure under the reference law makes the
reference conditional law in every slot a fresh iid block.  Discarding
the nonnegative entropy terms for the gaps, past, and other pieces
leaves the asserted upper bound for the slots.  This proof also
applies to finite restrictions of an infinite past, followed by
monotone increase of the observed sigma fields.

For an independent grid displacement $U$, apply the argument
conditionally on $U$ and integrate.  The original word law is
independent of $U$, and the selected law before grid sampling has
the same bounded-density property.  We do not condition on a
sample-dependent favorable displacement before this calculation.
\end{proof}

\begin{corollary}[Central oscillations at intermediate levels]
\label{arm:central-oscillations}
For a projected address code, assign to every occupied node $v$ an
interval of length at most $A r_v$ centered within its location cell,
where $A$ is fixed and $r_v$ is its viewing radius.  Centers rounded
using the terminal address are also allowed, provided they lie in the
specified cell.  Write
\[
 S_v=\frac{\operatorname{osc}_{I_v}Z}{\sqrt{r_v}}.
\]
Suppose $\mathbb E_{\mathbb Q}\sum_jm_j=O(N)$, the expected code length
is $O(N)$, and the selected density bound is $e^{O(N)}$.  Then
\[
 \mathbb E_{\mathbb Q}\sum_v S_v=O(N).
\]
In particular, the sum of the central oscillation amplitudes of the
occupied lines has bounded first moment at a uniform intermediate
viewing level.  The constants may depend on the fixed grid ratios,
padding, and $A$, but not on the exploration-list cutoff.
\end{corollary}

\begin{proof}
Replace each $I_v$ by the deterministic interval
\[
 \widehat I_v=
 [\inf v-A r_v/2,\ \sup v+A r_v/2],
\]
where $v$ denotes the code's location cell.  This interval contains
$I_v$ regardless of the chosen in-cell center.  Its length is exactly
$(A+R^{-2})r_v$.  At a fixed level, location cells have common length
$r_v/R^2$, so these enclosures have overlap multiplicity at most
$C_A R^2$.  Across levels their lengths are geometric.  The proof of
Lemma~\ref{arm:oscillation-sum} gives constants $\theta,c>0$ such
that, for every fixed code,
\[
 \mathbb E_{\mathbb P}
 \exp\left\{\theta\sum_v S_v-c\sum_jm_j\right\}\le1.
\]
The same holds under the product reference in which the code has
its Kraft weights and is independent of the raw word.  The entropy
variational inequality and Lemma~\ref{arm:entropy-allocation} imply
\[
 \theta\mathbb E_{\mathbb Q}\sum_vS_v
 \le c\mathbb E_{\mathbb Q}\sum_jm_j+
       \kappa+\mathbb E_{\mathbb Q}\ell(C)=O(N).
\]
Division by the number of intermediate levels proves the last claim.
\end{proof}

\begin{remark}[Arbitrary off-center bands at a coarsened level]
\label{arm:coarsened-bands}
The entropy allocation is not restricted to seed probe levels.  Fix
finitely many bands $a r\le |t-t_v|\le A r$, with $0<a<A<\infty$,
at the intermediate viewing radius $r$.  Refine the location cells
so their length is at most $a r/8$, retain only a residue class of
levels with successive radius ratio greater than $8A/a$, and increase
the ancestor padding in Lemma~\ref{arm:regular-slots}.  The same
ancestor/descendant proof then makes the enlarged test bands pairwise
disjoint: descendants lie in the inner hole, and incomparable nodes
lie outside the padded ancestor.  The residue class is chosen uniformly
and independently, so averaging over it preserves the uniform-level
marginal.  Its spacing is fixed once the finite band request is fixed.

For origins rounded from terminal address cells, slow coarsening makes
the rounding error $o(r)$.  The deterministic enlarged bands therefore
contain every actual requested block, for large enough coarsening.
One applies Lemma~\ref{arm:entropy-allocation} to the enlarged raw
slots and then takes the requested observations as measurable
functions there.  Structural and boundary deletions are bounded by
the same marker and fractional-occupancy errors.  A continuous
logarithmic phase, if used, is integrated as independent reference
side information; equivalently, stationarity under the fixed grid
dilation already suffices for the zero-or-infinite conclusions about
pinned heights and first sign-change times.
\end{remark}
